\section{Introduction}

Much operational computer work is \emph{repeated}: the business intent is stable,
the sequence is mostly known, errors are costly, and the same task may run
hundreds of times. Reasoning through such a workflow on every run is both
wasteful and unsafe. Wasteful, because a general model re-derives the same plan
each time, paying latency, inference cost, and run-to-run variation. Unsafe,
because the usual success signal is the \emph{screen}: a model or a script sees a
rendered ``Saved'' banner and stops---but a banner is not a persisted write. The
same green pixels can follow a partial save, a duplicate submission, a stale
overwrite, or a silently rejected transaction. The failure mode that matters in
production is not the visible crash; it is the \emph{silent wrong action} that
reports success.

Graphical user interfaces remain the practical integration boundary for legacy,
third-party, and remote applications, so this problem cannot be assumed away with
an API. Selector-based automation is fast and deterministic when stable structure
is available but brittle under interface drift, while visual scripting extends to
pixel surfaces \cite{sikuli}. Recent computer-use agents instead observe the
screen and choose each action with a general model
\cite{anthropic2024computeruse,uitars2025,cheng2024seeclick}, a useful
formulation for previously unseen tasks and broad benchmarks
\cite{webarena,workarena,osworld}. Neither line closes the safety gap: RPA does
not check the system of record, and an agent re-reasons every run yet still trusts
the screen.

\system{} treats a successful human demonstration as source material for a
\emph{program} rather than as context for a model on every run
(Figure~\ref{fig:pipeline}). Compilation extracts parameters, redundant target
evidence, postconditions, and risk metadata. Replay follows the compiled program
deterministically and makes zero model calls on the healthy path; models are
optional repair tiers, not the controller. Where a step is consequential, an
independent verifier checks the effect against the application's own
system of record, and the runtime refuses rather than guessing when it cannot
verify. The distinction the paper foregrounds is not whether a task
\emph{completed} but whether identity and external effects were \emph{verified},
and whether the system can \emph{refuse} when evidence is insufficient. In an
injected-fault study this distinction is decisive, and we measure it through the
live replayer's governed actuation and verification path and judge it by an
independent system-of-record ground
truth: screen-only verification silently accepted 54 of 90 fault runs, an
out-of-band effect check cut that to 9 of 90, and a complete system-of-record
read path to 0 of 90 (Section~\ref{sec:results}). Conditioned on the 72 runs in
which a wrong effect actually occurred, that is 75.0\%, then 12.5\%, then zero. The same gap replicates on a
second, non-healthcare lending domain and on two multi-application workflows
spanning email, spreadsheets, documents, and two systems of record.

This paper makes five contributions:
\begin{enumerate}
  \item A demonstration compiler and backend-neutral workflow representation
        for deterministic repeated GUI execution.
  \item A resolution ladder that prefers structured identity, then local visual
        evidence, and records bounded drift repairs as reviewable changes.
  \item System-of-record effect verification that treats the screen as a weak
        oracle and the application's own state as the strong oracle, closing the
        gap between a rendered success and a persisted effect.
  \item Governance mechanisms that distinguish runnable from certified, verify
        identity and system effects where configured, and refuse rather than
        silently degrade. Three are load-bearing for the thesis and are
        described in Section~\ref{sec:governance}: an explicit \emph{transaction
        outcome} taxonomy backed by a privacy-safe effect journal and an
        at-most-once idempotency ledger, so an uncertain write is reconciled
        rather than retried; a \emph{governed repair promotion lifecycle} in
        which a runtime repair is only ever a candidate until human review,
        contract-invariant enforcement, replay and fault campaigns, a
        digest-bound approval, and a bounded auto-rollback canary have all
        passed; and \emph{explicit surface selection}, which refuses an implicit
        browser default in production and refuses to run a workflow qualified on
        one execution surface on another.
  \item An evaluation protocol that reports task success \emph{together with}
        silent incorrect success and over-halt, with a failure taxonomy and
        \emph{machine-checked paper constants} that bind every headline number
        to a released raw artifact or bounded aggregate summary. The protocol is
        instantiated as EffectBench, a standalone versioned benchmark, and is
        exercised in a second non-healthcare domain and on two multi-application
        workflows so that the metric is not tied to one fixture. Raw grown
        corpora, deployment-derived tuning, and target-specific recipes remain
        outside the public release. The machine-check establishes transcription
        fidelity---that the prose faithfully cites the released evidence---not
        the validity of the underlying measurement; a sound number and an
        unsound one are bound alike.
\end{enumerate}

The evaluation deliberately separates a broad product architecture from the
scope of each measured result. Browser workflows supply the comparative and
breadth studies. Additional task-bound qualifications cover Windows UIA,
native macOS, and real-network RDP with independent effect oracles and ambiguity
or stale-target refusals. Two multi-application benchmarks add branching,
exception paths, and cross-system effects. These results establish working
substrate and workflow mechanisms
for the named environments; they are not a population estimate for arbitrary
enterprise applications. No external organization has deployed the system, so
nothing here is field-reliability evidence.
Citrix ICA/HDX is a pixel-only remote-display
environment in the same class as the measured RDP backend, and we report a
qualification of the shipping Citrix Workspace backend against a
\emph{deterministic synthetic stand-in} that reproduces ICA/HDX conditions in
process. That stand-in is not Citrix: it exercises no HDX codec, no ICA
compression, and no Workspace-client transport, and its released artifact
records \texttt{is\_real\_ica\_hdx: false}. The measured remote-display result
in this paper is RDP over Aardwolf into a Parallels Windows guest, and RDP
evidence does not transfer to Citrix either.
The limitations section states this scope precisely and names the
specific validation it would take.
